\chapter{Andreev reflection and subgap states}\label{ch:background}

\section{Three devices, one superconducting environment}

A superconducting weak link is a small region through which two superconductors
exchange electrons and Cooper pairs. Its microscopic structure matters. A clean
constriction is described by transmission channels, a confined quantum dot by
orbitals and charging energies, and a magnetic adatom by its exchange coupling to
the host. All can bind excitations below the superconducting gap, but the mechanism
and the appropriate approximations differ. \Cref{fig:primer-devices} compares these
devices before we introduce their energy levels.

In a \emph{quantum point contact} (QPC), transverse confinement leaves a few
conducting channels. A channel has a normal-state transmission probability
$0\le\tau_n\le1$ and is defined as an eigenmode of the scattering problem.
An atomic contact is a particularly short constriction, and even a single
atom can support several channels through its valence orbitals. A gate-defined
semiconductor QPC allows the channels to be tuned electrostatically. A
\emph{quantum dot} instead confines charge sufficiently strongly that the energy
cost of adding a second electron must be included explicitly. Its contact
transmissions generally depend on energy, occupation and correlations. The
distinction is determined by the physical regime and applies across different
device materials.

\begin{figure}[htb]
\centering\manualfigure{primer_devices}
\caption{Three routes to localized superconducting excitations: a short
constriction between superconductors, a gate-controlled dot with tunnel barriers,
and a magnetic adatom on a superconducting surface. The arrows identify control
or probe couplings; the device sizes and shapes are schematic.}
\label{fig:primer-devices}
\end{figure}

The phase difference $\phi=\phi_R-\phi_L$ controls interference between pair
amplitudes. A gate changes the dot energy and often also the tunnel barriers. A
magnetic field can split a spin doublet and, in a real device, also affect the
superconducting leads. The principal energies are the gap $\Delta$, charging
energy $U$, total hybridization $\Gamma$, detuning $\eps$, Zeeman energy $B$ and
thermal energy $k_BT$. Ratios of these scales determine which physical picture is
useful. \Cref{ch:model} defines the model parameters precisely; the present
chapter develops the non-interacting reference picture.

\section{What is inside the gap?}\label{sec:gap-picture}

Recall that a BCS condensate is a coherent paired state, whereas a Bogoliubov
quasiparticle is an excitation above it. For a homogeneous $s$-wave superconductor,
$E_k=\sqrt{\xi_k^2+\Delta^2}$, so propagating single-quasiparticle excitations
begin at $\Delta$. For a constant normal-state density of states $\rho_0$
per spin, the ideal bulk electron density of states is
\begin{equation}\label{eq:primer-dos}
 \rho_{\mathrm{BCS}}(E)=\rho_0\frac{|E|}{\sqrt{E^2-\Delta^2}}
 \Theta(|E|-\Delta).
\end{equation}
Here $E$ is measured from the chemical potential and can be positive or negative
because an electron probe can add or remove an electron. The positive
quasiparticle excitation energy must be distinguished from this signed probe
energy.

A local defect or weak link changes the Hamiltonian and can produce discrete
spectral poles inside this bulk gap. Their wavefunctions decay into the leads;
they are \emph{bound states}. Thus an ideal hard bulk gap and local subgap peaks
are compatible. A soft gap instead describes appreciable spectral background
inside the nominal gap, which can have several microscopic origins. A population
of nonequilibrium quasiparticles specifies the occupations of the available states.

\begin{figure}[htb]
\centering\manualfigure{primer_gap}
\caption{(a) Ideal, unbroadened BCS density of states, with the gap edges marked.
The plotting range truncates the divergences.
(b) Schematic local spectral poles of a bound state, with unequal electron-addition
and removal weights, represented by the pole heights.
(c) Illustrative decay envelopes: a state nearer the gap edge
extends farther into the superconductor.}
\label{fig:primer-gap}
\end{figure}

For a clean lead in the Andreev approximation, an envelope has the characteristic
length
\begin{equation}
 \xi_E=\frac{\hbar v_F}{\sqrt{\Delta^2-E^2}},\qquad |E|<\Delta,
\end{equation}
where $v_F$ is the Fermi velocity. Oscillations on the Fermi-wavelength scale and
geometric prefactors are omitted in this envelope picture. A level approaching
the gap edge spreads into the reservoir and eventually merges with its continuum.
Thus a subgap wavefunction generally extends over both dot and leads. Its dot spectral
weight can change substantially while its total
normalization remains one.

\begin{physics}[A gap is a statement about excitations]
The condensate can carry a stationary supercurrent without creating a population
of quasiparticles. A single bulk quasiparticle costs at least $\Delta$; creating
two from an even-parity vacuum costs at least $2\Delta$. Localized states change
the corresponding thresholds. Whenever a diagram shows a shaded continuum,
ask which final states and which parity sector it represents.
\end{physics}

\section{Andreev reflection: exchanging a pair with the condensate}

Consider an electron incident from a normal region on an ideal transparent
normal--superconductor interface, with excitation energy below $\Delta$, the
threshold for bulk quasiparticle propagation. It can be reflected as a
hole while a Cooper pair is transferred into the condensate. For singlet pairing,
the electron and the missing electron represented by the hole have opposite spin.
The reflected hole approximately retraces the incident path when the Fermi energy
is much larger than $\Delta$ and the excitation energy. This is
\emph{Andreev reflection}. A barrier also permits ordinary electron reflection;
the relative amplitudes are part of the scattering problem.

\begin{figure}[htb]
\centering\manualfigure{primer_reflection}
\caption{Andreev conversion at one interface and an electron--hole round trip
between two superconductors. Electron and hole arrows denote group-velocity
directions of coherent amplitudes. Pair transfer accompanies conversion; the
trajectory describes the phase quantization of coherent electron--hole amplitudes.
The short-junction limit neglects the energy-dependent phase
accumulated while traversing the normal region.}
\label{fig:primer-reflection}
\end{figure}

Two interfaces can confine a coherent electron--hole orbit. Each conversion adds
an energy-dependent phase involving $\arccos(E/\Delta)$ and the phase of the
condensate. For an ideally transmitting short link, a round trip requires
$2\arccos(E/\Delta)=\pm\phi+2\pi m$ for an integer $m$, with the sign distinguishing
the two propagation directions. This produces a discrete phase-dependent level.
Normal reflection mixes the trajectories. The resulting expression for one
spin-degenerate channel is~\cite{beenakker1991josephson,beenakker1991universal}
\begin{equation}\label{eq:contact-abs}
 E_A(\phi)=\Delta\sqrt{1-\tau\sin^2(\phi/2)},\qquad E_A\ge0.
\end{equation}
The two leads have equal conventional gaps, the normal scattering matrix is
effectively energy independent over that gap, and Coulomb interaction is omitted.
The short-junction condition is a dwell-time condition,
$\Delta t_{\mathrm{dwell}}/\hbar\ll1$. Assess it through the dwell time: even a
geometrically small dot can have a narrow resonance and a long dwell time.

At zero phase the level reaches the gap edge. Near $\phi=\pi$, increasing
transmission lowers its energy; for $\tau<1$ a finite minimum remains. The
zero-energy touching at perfect transmission is sensitive to reflection and
occupation. Establishing topological protection or a protected
$4\pi$-periodic equilibrium Josephson effect requires additional evidence.

\section{One Andreev orbital, four many-body states}\label{sec:abs-occupations}

The Bogoliubov--de Gennes (BdG) equation is a one-particle eigenvalue problem in an
enlarged electron--hole space. Its $+E_A$ and $-E_A$ partners describe a
particle--hole redundancy of the same physical excitation.
After choosing the positive-energy modes, the many-body Hamiltonian of a
spin-degenerate Andreev orbital can be written
\begin{equation}\label{eq:andreev-occupations}
 H_A=E_A\left(\gamma_\up\dg\gamma_\up+
                    \gamma_\dn\dg\gamma_\dn-1\right).
\end{equation}
The four Fock states are
\begin{align}
 \ket g&=\ket{0_A}, & E_g&=-E_A, & P_g&=+1,\nonumber\\
 \ket{o_\sigma}&=\gamma_\sigma\dg\ket{0_A}, & E_{o_\sigma}&=0, & P_{o_\sigma}&=-1,\nonumber\\
 \ket e&=\gamma_\up\dg\gamma_\dn\dg\ket{0_A}, & E_e&=+E_A, & P_e&=+1.
\end{align}
The notation $0_A$ counts \emph{bound quasiparticles}, while $\ket0$ denotes the empty
electron dot. The state $\ket{0_A}$ already includes superconducting
pairing and all filled negative-energy electronic states of the reference vacuum.

\begin{figure}[htb]
\centering\manualfigure{primer_occupations}
\caption{Four occupations of a spin-degenerate Andreev orbital. Single-electron
addition or removal connects even and odd parity at excitation energy $E_A$.
A pair transition connects the two even states at $2E_A$. The odd states are
degenerate here; their splitting in a spin qubit requires additional physics.
The arrows indicate transitions allowed by parity; their strengths depend on the
matrix elements of the chosen drive.}
\label{fig:primer-occupations}
\end{figure}

This elementary example resolves a common apparent discrepancy between
experiments. A weak normal tunnel probe can produce resonances at biases
$eV=\pm E_A$, whereas a parity-preserving microwave probe can absorb at
$hf=2E_A$. The two probes access different transitions of the same system. The
interaction in a quantum dot changes the level spacings and wavefunctions.
Calculate each transition from the appropriate many-body energy difference
and matrix element of the physical probe operator.

Parity is conserved because pairing changes electron number by two. It remains
well defined when particle number itself is not. A state of definite parity can
have a non-integer \emph{mean local charge}; its parity follows from the occupied
Fock configurations. For a coupled dot and leads, it is the parity of the combined
system that labels an eigenstate. Tunnelling produces fluctuations of the dot charge within that state.

\section{Why a bound state carries supercurrent}

An energy level that depends on phase contributes a current. For a nondegenerate
stationary state, $I=(2e/\hbar)\partial_\phi E$. Applying this relation to the
even ground state of \cref{eq:andreev-occupations} gives
\begin{equation}\label{eq:contact-current}
 I_g(\phi)=\frac{e\Delta}{2\hbar}
 \frac{\tau\sin\phi}{\sqrt{1-\tau\sin^2(\phi/2)}}.
\end{equation}
This uses the manual's phase and positive-current convention. The excited even
state carries $-I_g$ and either odd state carries zero current in this ideal
short-channel model. The two quasiparticles cancel the ground-state contribution
one by one. For equal gaps and energy-independent short-junction scattering, the
continuum contribution to the equilibrium energy is phase independent. For general
dots, long junctions or unequal gaps, include the continuum contribution explicitly.

\begin{figure}[htb]
\centering\manualfigure{primer_contact}
\caption{Analytic short-channel spectrum and ground-state current for the
transmissions indicated. Energy is normalized by $\Delta$ and current by
$e\Delta/\hbar$. Higher transmission produces a deeper bound state and a more
skewed current--phase relation. These curves assume the even ground state at zero
temperature and use \cref{eq:contact-abs,eq:contact-current}.}
\label{fig:primer-contact}
\end{figure}

At small transmission the current is nearly sinusoidal. At high transmission
the energy changes rapidly near $\pi$, giving a strongly nonsinusoidal current.
An excitation changes the occupied many-body branch and hence the current; this
provides one method of Andreev-state readout. The current depends on the full
phase-dependent energy reference. Differentiating a plotted excitation
$E_m-E_g$ gives the current difference $I_m-I_g$.

If all four occupations equilibrate with a reservoir, their partition function is
$Z=\ee^{\beta E_A}+2+\ee^{-\beta E_A}$, where $\beta=(k_BT)^{-1}$. The resulting
equilibrium current is $I_g\tanh[E_A/(2k_BT)]$. A device trapped in one parity
sector samples a restricted ensemble, whose current is obtained by averaging over
that sector. Quasiparticle poisoning is the process that switches
the bound subsystem between such occupation sectors.

\begin{physics}[Read the figures]
Suppose a microwave line follows twice the energy of a tunnel-spectroscopy peak.
Which two occupations are connected in each experiment? The microwave line can
connect $\ket g$ to $\ket e$; tunnelling connects $\ket g$ to $\ket{o_\sigma}$.
If one quasiparticle is trapped, the pair transition from $\ket g$ loses its
initial population, so the line disappears while the bound level can remain
inside the gap. The same trapped odd state carries zero current in the ideal
model; an interacting dot doublet can carry a finite current.
\end{physics}

\paragraph{Reading further.}
The scattering picture and its connection to quantum dots are reviewed by
Mart\'in-Rodero and Levy Yeyati~\cite{martinrodero2011josephson}. Aguado's
perspective~\cite{aguado2020perspective} connects the occupation picture to
semiconductor superconducting qubits. Glazman and Catelani's lecture
notes~\cite{glazman2021bogoliubov} develop the distinction between the condensate,
quasiparticles and circuit degrees of freedom. The following chapter introduces
the charging energy that makes the dot a genuinely interacting artificial atom.
